\input{estilo}
\begin{document}
\begin{figura}\begin{tikzpicture}
  % las dos mitades con la misma grilla: memorias arriba, CPUs abajo, 3 cm entre CPUs
  % multiprocesador: los dos CPUs llegan directo a la misma memoria
  \node[caja, minimum width=40mm, minimum height=10mm] (mem) at (1.5,0) {memoria};
  \foreach \x/\n in {0/a, 3/b} {
    \node[caja] (cpu\n) at (\x,-2) {CPU};
    \draw (cpu\n) -- (\x,0 |- mem.south);
  }
  \node[font=\small] at (1.5,-3.95) {multiprocesador};
  % sistema distribuido: cada CPU con su memoria, unidos solo por la red
  \foreach \x/\n in {6.6/c, 9.6/d} {
    \node[caja] (m\n) at (\x,0) {M};
    \node[caja] (cpu\n) at (\x,-2) {CPU};
    \draw (m\n) -- (cpu\n);
    \node[marco, dashed, draw=suave, fit=(m\n)(cpu\n), inner sep=4mm] {};
  }
  \draw (cpuc) -- (cpud);
  % lo que no existe: un CPU que escriba en la memoria del otro
  \draw[dashed, draw=acento, shorten >=2pt] (cpuc.north east) -- (md.south west);
  \pic[scale=0.6] at ($(cpuc.north east)!0.5!(md.south west)$) {rechazo};
  \draw[llave, decoration={mirror}] (7.55,-2.95) -- (8.65,-2.95)
    node[midway, below=7pt, etiq, text=suave] {enlace de red};
  \node[font=\small] at (8.1,-3.95) {sistema distribuido};
\end{tikzpicture}\end{figura}
\end{document}
